%====================================================
% Mathematics frames continue under Methodology (no new \section)

\begin{frame}{Tokenisation and \texttt{[CLS]}}
	\begin{itemize}
		\item WordPiece splits English complaint \(X\) into subword tokens (max length \(M \leq 128\)).
		\item Special tokens: \texttt{[CLS]} at position \(0\), \texttt{[SEP]} at the end, \texttt{<PAD>} as needed.
		\item \textbf{Why:} a fixed-length token sequence is the only input BioBERT can read.
	\end{itemize}
	\begin{block}{Summary vector}
		After \(L=12\) encoder layers, the first row is the complaint summary:
		\[
		  h_{\mathrm{[CLS]}} \in \mathbb{R}^{768}.
		\]
		Classification reads this vector only, not every token separately.
	\end{block}
\end{frame}

\begin{frame}{Worked Example: Token IDs}
	{\scriptsize Complaint (opening clause in table): \textit{``I have a headache and feel feverish.''}; then pad to \(M=128\).}
	\vspace{0.1em}
	\begin{center}
	\scriptsize
	\begin{tabular}{|c|l|l|r|}
	\hline
	\(i\) & Raw & Token & ID \\
	\hline
	0 & n/a & \texttt{[CLS]} & 101 \\
	1 & I & i & 1045 \\
	2 & have & have & 2031 \\
	3 & a & a & 1037 \\
	4--5 & headache & head+\texttt{\#\#ache} & 7994, 8772 \\
	6 & and & and & 1998 \\
	7 & feel & feel & 2514 \\
	8--9 & feverish & fever+\texttt{\#\#ish} & 9643, 6804 \\
	10 & . & . & 1012 \\
	11 & n/a & \texttt{[SEP]} & 102 \\
	12--127 & n/a & \texttt{<PAD>} & 0 \\
	\hline
	\end{tabular}
	\end{center}
	\vspace{0.05em}
	\begin{itemize}\scriptsize
		\item \texttt{\#\#} = continuation of a word (not a new word), so ``headache'' uses two positions.
		\item \texttt{[CLS]} starts the sequence (later the 768-d summary); \texttt{[SEP]} ends real text; \texttt{<PAD>} fills to 128.
		\item BioBERT only sees this list of IDs, never the raw English string.
	\end{itemize}
\end{frame}

\begin{frame}{Input Embeddings \(H^{(0)}\)}
	\begin{block}{Three embeddings summed}
		For token \(t_i\) at position \(i\) (segment \(s=0\)):
		\[
		  E(t_i)_d = E_{\mathrm{word}}(t_i)_d + E_{\mathrm{pos}}(i)_d + E_{\mathrm{seg}}(0)_d,
		  \quad d=1,\ldots,768.
		\]
	\end{block}
	\begin{itemize}
		\item Stack rows to form the input matrix \(H^{(0)} \in \mathbb{R}^{M \times 768}\).
		\item BioBERT embedding space groups related clinical terms after PubMed pre-training \citep{lee2020biobert}.
		\item \textbf{Why:} meaning + order + segment identity enter the model before attention runs.
	\end{itemize}
\end{frame}

\begin{frame}{Attention in One Minute}
	\begin{block}{What attention does}
		Each word decides how much to \emph{listen} to every other word
		(e.g.\ ``no'' should weaken ``pain''; ``severe'' should strengthen ``headache'').
	\end{block}
	\begin{itemize}
		\item \(Q\): what am I looking for?
		\item \(K\): what do I offer as a match?
		\item \(V\): what content do I pass if chosen?
		\item Twelve heads of size \(64\): \(12 \times 64 = 768\) (full BioBERT width).
	\end{itemize}
	\begin{exampleblock}{Next slide}
		The formula is the same idea in math: score matches, scale, softmax weights, mix values.
	\end{exampleblock}
\end{frame}

\begin{frame}{Self-Attention (Encoder Layers)}
	\begin{block}{Scaled dot-product attention}
		\[
		  \mathrm{Attention}(Q,K,V)
		  = \mathrm{softmax}\!\left(\frac{QK^{\top}}{\sqrt{d_k}}\right)V,
		  \qquad d_k = 64.
		\]
	\end{block}
	\begin{itemize}
		\item \(Q=HW_Q\), \(K=HW_K\), \(V=HW_V\) from the current hidden matrix \(H\).
		\item Twelve heads in parallel; outputs concatenated and projected back to 768-d.
		\item Stack \(L=12\) identical encoder layers with residual connections \citep{vaswani2017attention}.
		\item \textbf{Why:} each symptom token can weigh every other token (e.g.\ negation + pain).
	\end{itemize}
\end{frame}

\begin{frame}{Softmax Classification}
	\begin{block}{Five-class head}
		\[
		\begin{aligned}
		  \mathbf{z} &= W h_{\mathrm{[CLS]}} + \mathbf{b}, \\
		  \hat{y}_c &= \frac{\exp(z_c)}{\sum_{j=1}^{5}\exp(z_j)}.
		\end{aligned}
		\]
		Each \(\hat{y}_c\) is the model's belief that the complaint is acuity level \(c\)
		(levels \(1\)--\(5\)); the five beliefs add up to \(1\).
	\end{block}
	\begin{itemize}
		\item \textbf{Chosen level \(\hat{c}\):} pick the level with the \emph{largest} probability
		(e.g.\ if Yellow has \(0.72\), choose Yellow).
		\item \(W \in \mathbb{R}^{5 \times 768}\), \(\mathbf{b} \in \mathbb{R}^{5}\).
		\item \textbf{Why:} turns the 768-d summary into five acuity scores that sum to one.
	\end{itemize}
\end{frame}

\begin{frame}{Colour Map \(f_{\mathrm{SATS}}\) and Pathway \(P(C)\)}
	\begin{block}{Acuity \(\to\) SATS colour}
		\[
		  f_{\mathrm{SATS}}(\hat{c}) =
		  \begin{cases}
		    \mathrm{Red} & \hat{c}=1 \\
		    \mathrm{Orange} & \hat{c}=2 \\
		    \mathrm{Yellow} & \hat{c}=3 \\
		    \mathrm{Green} & \hat{c}\in\{4,5\}
		  \end{cases}
		\]
	\end{block}
	\begin{exampleblock}{Pathway card}
		\(P(C) = \bigl(T_{\max},\;\mathrm{destination},\;\mathrm{actions},\;\mathrm{escalation}\bigr)\)
		aligned to \examplehospital{} protocols.\\
		\textbf{Why:} routing uses colours and destinations, not raw acuity integers.
	\end{exampleblock}
\end{frame}

\begin{frame}{Clinical Gate and End-to-End Chain}
	\begin{block}{Relevance gate}
		\[
		  g(X) = \mathrm{clip}\!\Bigl(\sum_j s_j(X) - p(X),\, 0,\, 1\Bigr),
		  \qquad \text{proceed iff } g(X) \geq 0.35.
		\]
		Symptom lexicon scores \(s_j\); greeting / non-clinical penalty \(p(X)\).
	\end{block}
	\begin{alertblock}{Full path}
		\[
		  X \xrightarrow{g} \tau(X) \to H^{(0)}
		  \xrightarrow{12} h_{\mathrm{[CLS]}}
		  \to \hat{\mathbf{y}}
		  \xrightarrow{f_{\mathrm{SATS}}} C_{\mathrm{NLP}}
		  \to P(C_{\mathrm{NLP}}).
		\]
		Skip BioBERT on greetings; one chain from text to pathway card.
	\end{alertblock}
\end{frame}
